% id: 124-17
% book: 베이직쎈 확률과 통계 · 07 통계적 추정
% section: 개념 38 표본평균의 평균, 분산, 표준편차
% figure: none
% answer: 
% answer_source: 
% variant_level: 0 (원본 전사)
% 이 파일은 build.mjs 가 items.json 에서 만든다. 고칠 때는 items.json 을 고친다.
\smash{\raisebox{1.5mm}{\dmkichul{베이직쎈 확률과 통계 124쪽 17번}}}
모집단 1, 3, 5에서 크기가 2인 표본을 임의추출할 때, 다음을 구하시오.
\dmsubprob{1} 표본평균 $\overline{X}$의 확률분포를 나타낸 표\par\smallskip\centerline{{\setlength{\tabcolsep}{5pt}\renewcommand{\arraystretch}{2.2}\begin{tabular}{|c|c|c|c|c|c|c|}\hline $\overline{X}$ & \makebox[2.2em]{} & \makebox[2.2em]{} & \makebox[2.2em]{} & \makebox[2.2em]{} & \makebox[2.2em]{} & 합계 \\ \hline $\mathrm{P}(\overline{X}=\overline{x})$ & & & & & & \\ \hline\end{tabular}}}\par\smallskip
\dmsubprob{2} $\mathrm{E}(\overline{X})$, $\mathrm{V}(\overline{X})$, $\sigma(\overline{X})$
